\subsection{Output bounds and termination}
\label{alg:termination-subsection}

Before discussing accuracy, we prove that every execution is finite and
that all recursive inputs and outputs have their stated types. These
facts must survive failed comparisons: a failed tape may change the
answers, but it must not invalidate the program or its resource bounds.
Recall that $d(R)$ is the least nonnegative integer satisfying
$(9/10)^{d(R)}R<1$.

\begin{proposition}\label{prop:termination}
For every permitted $\Solve$ input, the execution and all its subroutine
executions terminate on every random tape. The outputs of $\Solve$
and $\Ask$ are integer vectors in their respective input boxes.
The outputs of $\Orient$ and $\OrientPlus$ have the stated choice
and list types, and each listed set has input measure at most $1/2$.
Moreover:
\begin{enumerate}[label=(\roman*)]
\item Tightening makes at most $8ns$ updates and $16ns+2$ calls to
$\Ask$.
\item A chain of main $\Solve$ continuations contains at most
\begin{equation}\label{eq:main-chain-length}
 K=8s(g+1)+1
\end{equation}
solver invocations.
\item One $\Orient$ trial visits at most $g+1$ orientation invocations
and returns at most $2g+2$ sets. Consequently $\OrientPlus$ returns
at most $J(2g+2)$ sets, counting repetitions.
\end{enumerate}
\end{proposition}

\begin{proof}
We induct on the solver rank and establish the output assertions at the
same time as termination. At rank zero, $R<1$, and the endpoint rule
returns immediately in the given box. Fix a positive rank and assume
these properties for every smaller-rank solver input.

An ordinary budget change multiplies $R$ by $4/5$. Reweighting a set
with $\mu(H)\le1/2$ multiplies it by
$(3/5)(1+\mu(H))\le9/10$. For either change,
\[
 (9/10)^{d(R)-1}R_{\mathrm{child}}
 \le(9/10)^{d(R)}R<1,
\]
so the child has smaller rank. Such $\Ask$ calls terminate by the
induction hypothesis. Their medians remain integer vectors in their
input boxes, irrespective of comparison accuracy.

\paragraph{Orientation at the current rank.}
The only solver calls within orientation are through budget-reduced
$\Ask$. At width $B\le128s$, $\Basic$ terminates by
Lemma~\ref{lem:box-solution}; the returned set is either $P$ of
measure at most $1/2$ or the complement of a set of measure greater
than $1/2$. At larger width, both
$D=\floor{B/(32s)}$ and $E_0=\floor{B/8}$ are positive. Each reduced
call uses a legitimate centered box with reference width
$2D\le B/(16s)$. Its result lies between $q_j-D$ and $q_j+D$.
The integer floor of their average lies between these same integer
endpoints. Since $q_j\in[a,b]$, subsequent clipping keeps the next
center in $[a,b]$ and cannot increase its distance from $q_j$. Therefore
\begin{equation}\label{alg:center-step-bound}
 |(q_{j+1})_i-(q_j)_i|\le D\qquad(i\in V).
\end{equation}
This proves that all $s$ loop steps terminate with valid centers.

The early returns again use complements of sets with measure greater
than $1/2$. When both tests fail, the appended $P,N$ each have measure
at most $1/2$. The child box has positive width $2E_0\le B/4$ and
unchanged support parameters, hence unchanged $\mu$. A separate
induction on this decreasing positive integer width proves orientation
termination and the measure assertion. After $r$ child edges its width
is at most $B_*/4^r$. Since $g>\log_2 B_*$, there are at most $g$
edges, or $g+1$ invocations. The final one returns one set and each
ancestor adds two, giving $1+2g\le2g+2$ entries. Repeating this finite
trial $J$ times gives a well-defined majority because $J$ is odd, and
concatenation proves the list bound for $\OrientPlus$.

\paragraph{Tightening at the current rank.}
The solver's small-width branch is already covered by basic iteration.
Otherwise $B>128s$, and
\begin{equation}\label{alg:delta-bound}
 \delta=\floor{B/(4s)}\ge\frac{B}{8s}>0,
 \qquad0<C=B-\delta<B.
\end{equation}
Every tightening output is inside the current box. A lower update only
raises $a$, an upper update only lowers $b$, and both preserve $a\le b$.
At least one coordinate moves by $\delta$. The integer quantity
$\sum_i(b_i-a_i)$ starts at most $nB$ and decreases by at least
$\delta$ at each update. There are at most $nB/\delta\le8ns$ updates.
A pass contains at most two $\Ask$ calls and either restarts after an
update or is the final pass. This gives $2(8ns+1)=16ns+2$ calls.

\paragraph{The main continuation and its returned vector.}
Orientation now terminates by the argument just given. All of its sets
have the required measure on every tape, so every reweighted call has
smaller rank and terminates. The sole main successor has a valid subbox
and reference width $C$. Along a chain of such successors,
\[
 B_{r+1}\le B_r\left(1-\frac1{8s}\right),\qquad
 B_r\le B_*\exp\left(-\frac r{8s}\right).
\]
After $8s(g+1)$ edges this last bound would be below one, because
$g+1>\ln B_*$. In fact a chain stops already at width at most $128s$.
Thus it has at most $K$ invocations. This reasoning uses the reference
width, so it remains valid when a cap leaves the actual box unchanged.

Work backwards along this finite chain. All reweighted outputs lie in
the node's final tightening box, and the main output lies in its
specified subbox. Their coordinatewise maximum or minimum is still in
the tightening box, itself contained in the node's input box. Backward
induction therefore proves the solver output assertion and termination
at the current rank. No correctness event has been used at any stage.
Since orientation is invoked only after the test $R<1$ fails, all uses
of positive rank above apply to the actual calls. This completes the
rank induction.
\end{proof}
